\phantomsection
\section*{七、三类算法的协同优化框架}
\addcontentsline{toc}{section}{七、三类算法的协同优化框架}

前文分别从模型切分部署、业务感知 batching 和 LLM inference task offloading 三个角度讨论了 CEC-LLM 推理系统的关键算法。真实系统中，这三类算法并不是彼此独立的模块：模型切分部署决定模型组件、服务实例和推理阶段的空间分布；业务感知 batching 决定请求、prefill chunk 和 decode token 在实例内部或不同阶段资源池中的时间组织；task offloading 决定完整请求、复杂 workflow 以及 context/KV/runtime state 在端、边、云之间的路由、迁移和复用。三者共享模型位置、队列长度、KV cache 位置、链路状态、用户移动趋势和 SLO slack 等同一组系统状态，因此任何一类算法的局部优化都可能改变另外两类算法的可行域和收益。

本章不再重新划分第四、五、六章的所有单点方法，而是围绕“三类算法之间怎样协同”组织已有工作。具体而言，协同优化关注的不是某个模型如何单独切分、某个 batch size 如何单独选择，或某个请求如何单独卸载，而是这些决策在多时间尺度、多资源层级和多运行时状态下如何形成闭环。已有工作已经从局部角度揭示了这种耦合关系：model caching 会影响 request offloading 的可选路径，prefill/decode disaggregation 会改变 batching 和 KV handoff，KVCache-centric serving 会把状态位置纳入 request scheduling，移动场景下的 context migration 又会反过来影响长期 session routing。

本文将 CEC-LLM 推理中的三类算法协同关系划分为四类：
\begin{enumerate}
\item \textbf{Model / Service Placement -- Request Offloading 协同}：低频模型缓存、服务部署和高频请求路由之间的双时间尺度闭环。
\item \textbf{Model Partitioning -- Batching / Scheduling 协同}：模型切分、阶段解耦、continuous batching 和 serving scheduling 之间的容量匹配与时延控制。
\item \textbf{Request / Workflow Routing -- Runtime State Management 协同}：请求路径、agent workflow、KV cache、prefix cache、session state 和 context migration 之间的状态感知路由。
\item \textbf{Mobility-aware Hierarchical Control}：用户移动、负载热点变化、handover 和状态位置变化共同驱动的跨层闭环控制。
\end{enumerate}

这四类关系不是新的孤立算法类别，而是前文三类算法在系统运行时的交叉界面。第一类主要连接第四章的模型/服务部署和第六章的请求卸载；第二类连接第四章的模型切分和第五章的 batching；第三类连接第五章的状态感知调度和第六章的请求/workflow offloading；第四类则把移动 CEC 中的外部扰动纳入统一控制框架。表~\ref{tab:cec-collaborative-optimization-overview} 汇总了与这些协同关系直接相关的代表工作、主归属和方法关键词。

\begingroup
\scriptsize
\setlength{\tabcolsep}{1.4pt}
\renewcommand{\arraystretch}{1.35}
\emergencystretch=2em
\begin{longtable}{@{}>{\raggedright\arraybackslash\hspace{0pt}}p{0.035\textwidth} >{\raggedright\arraybackslash\hspace{0pt}}p{0.175\textwidth} >{\raggedright\arraybackslash\hspace{0pt}}p{0.095\textwidth} >{\raggedright\arraybackslash\hspace{0pt}}p{0.170\textwidth} >{\raggedright\arraybackslash\hspace{0pt}}p{0.155\textwidth} >{\raggedright\arraybackslash\hspace{0pt}}p{0.330\textwidth}@{}}
\caption{\textbf{CEC-LLM 三类算法协同优化相关工作总览}}
\label{tab:cec-collaborative-optimization-overview}\\
\toprule
\textbf{序号} & \textbf{代表工作} & \textbf{发表在哪} & \textbf{主归属} & \textbf{相关机制} & \textbf{方法关键词} \\
\midrule
1 & SLM-LLM-Collab \refcite{60} & TMC 2026 & Model / Service Placement -- Request Offloading & SLM caching；D2D / edge inference offloading & 长期缓存 SLM 或插件模型，短期在 local、D2D 和 edge LLM 之间选择卸载路径，体现模型缓存与请求卸载的双时间尺度协同 \\
\tabrowrule
2 & T2DRL \refcite{68} & TON 2026 & Model / Service Placement -- Request Offloading；Mobility-aware Control & long-context model caching；request management & 在 mobile edge 中联合管理 cached models、service requests 和资源匹配，使 long-context LLM agent 服务同时感知模型缓存、上下文长度和边缘资源 \\
\tabrowrule
3 & H-MEC-Placement \refcite{71} & TMC 2025 & Model / Service Placement -- Request Offloading & service placement；task scheduling；resource allocation & 面向 general MEC 说明服务部署、任务调度和资源分配在规划周期内存在强耦合，可作为 LLM 服务布局协同的基础参考 \\
\tabrowrule
4 & LVM-Offload \refcite{72} & INFOCOM 2026 & Model / Service Placement -- Request Offloading & foundation model selection；step offloading；resource allocation & 在 foundation model 服务中联合考虑模型选择、推理步数、任务卸载和资源分配，说明模型能力选择会改变后续卸载路径 \\
\tabrowrule
5 & DistServe \refcite{7} & OSDI 2024 & Model Partitioning -- Batching / Scheduling & prefill/decode disaggregation；phase-specific workers & 将 prefill 和 decode 分离到不同 worker pool，并分别优化并行策略、batching 和资源分配，收益取决于 KV handoff 与阶段队列协调 \\
\tabrowrule
6 & Splitwise \refcite{18} & ISCA 2024 & Model Partitioning -- Batching / Scheduling & prompt pool；token pool；KV transfer & 把 prompt computation 和 token generation 放到不同机器池，利用阶段资源画像差异，同时需要处理阶段间 KV transfer 和调度边界 \\
\tabrowrule
7 & Orca \refcite{4} & OSDI 2022 & Model Partitioning -- Batching / Scheduling & iteration-level scheduling；selective batching & 在 decode iteration 边界更新 active batch，说明 serving 调度粒度会直接影响模型实例利用率和请求等待时间 \\
\tabrowrule
8 & vLLM \refcite{5} & SOSP 2023 & Model Partitioning -- Batching / Scheduling；Runtime State Management & PagedAttention；KV block management & 通过 KV block 管理支撑 continuous batching，使 batching 容量不只取决于计算资源，也取决于 KV 显存碎片和状态管理能力 \\
\tabrowrule
9 & Llumnix \refcite{19} & OSDI 2024 & Model Partitioning -- Batching / Scheduling；Runtime State Management & live request migration；instance rescheduling & 在多个 LLM serving 实例之间迁移运行中请求及其状态，体现实例调度、batch 负载均衡和状态迁移之间的耦合 \\
\tabrowrule
10 & Mooncake \refcite{21} & USENIX FAST 2025 & Request / Workflow Routing -- Runtime State Management & distributed KV cache；global scheduling；PD disaggregation & 将 KV cache 作为全局状态层管理，调度器按 KV 本地性、节点负载和 SLO 选择 prefill/decode 路径 \\
\tabrowrule
11 & Preble \refcite{23} & ICLR 2025 & Request / Workflow Routing -- Runtime State Management & distributed prompt scheduling；prefix/KV locality & 把 prefix/KV cache 命中和负载均衡共同纳入请求分配，使 routing action 与后续 prefill 复用和 batch 组成相关 \\
\tabrowrule
12 & CachedAttention \refcite{24} & USENIX ATC 2024 & Request / Workflow Routing -- Runtime State Management & attention / prefix cache reuse & 说明 attention/prefix cache 的复用机会会改变请求服务路径和重复计算成本，是状态感知调度的重要依据 \\
\tabrowrule
13 & FuseSpill \refcite{66} & TPDS 2026 & Request / Workflow Routing -- Runtime State Management & KV spillover；auxiliary GPU scheduling & 显存受限时根据代价模型选择 KV cache spillover，并把溢出序列调度到辅助 GPU，改变原有 decode 执行位置 \\
\tabrowrule
14 & QoS-LLM-Router \refcite{55} & TMC 2025 & Request / Workflow Routing -- Runtime State Management & QoS-aware routing；queue / memory state & 将 edge LLM expert routing 建模为长期 QoS 优化，路由决策同时感知质量、响应长度、队列状态和 GPU memory \\
\tabrowrule
15 & CES-CB \refcite{56} & TSC 2025 & Request / Workflow Routing -- Runtime State Management & request cost prediction；edge batch selection & 利用 prompt embedding 和历史执行数据预测请求成本，在 edge/cloud 间选择请求集合，说明 routing 会影响队列阻塞和 batch 组成 \\
\tabrowrule
16 & MasRouter \refcite{62} & ACL 2025 & Request / Workflow Routing -- Runtime State Management & multi-agent routing；model / role selection & 根据 query 决定 multi-agent 协作模式、agent 数量、角色和底层 LLM；在 CEC 中需要进一步映射为 workflow offloading \\
\tabrowrule
17 & Conductor \refcite{63} & ICLR 2026 & Request / Workflow Routing -- Runtime State Management & natural-language workflow control & 训练控制器生成多模型 workflow、subtask 和上下文访问范围，提示 agent workflow 需要同时考虑模型选择和状态可见性 \\
\tabrowrule
18 & AFLOW \refcite{64} & ICLR 2025 & Request / Workflow Routing -- Runtime State Management & automatic agentic workflow search & 将 agentic workflow 表示为代码形式的 LLM 调用图，并用执行反馈搜索更优 workflow；在 CEC 中还需加入物理节点和状态路径 \\
\tabrowrule
19 & MGCO \refcite{61} & TSC 2026 & Mobility-aware Hierarchical Control & mobility-aware generative computation offloading & 用 Transformer MHA encoder 建模历史决策和未来移动状态，说明请求卸载不能只根据当前时隙状态短视决策 \\
\tabrowrule
20 & DT-MADTO \refcite{65} & TMC 2025 & Mobility-aware Hierarchical Control & digital twin；DAG offloading；mobility prediction & 通过 digital twin 预测移动和资源状态，联合决策 DAG 子任务卸载、带宽分配和中间数据等待时间 \\
\tabrowrule
21 & CA-AIGC-Migration \refcite{69} & TSC 2026 & Mobility-aware Hierarchical Control；Runtime State Management & context migration；mobility-aware service migration & 将完整模型迁移转化为 long context migration，在回复质量、AoI、计算延迟和传输成本之间做权衡 \\
\bottomrule
\end{longtable}
\endgroup


\phantomsection
\subsection*{7.1 三类算法的基本耦合关系}
\addcontentsline{toc}{subsection}{7.1 三类算法的基本耦合关系}

模型切分、batching 和 task offloading 共享同一组系统状态，包括模型副本位置、模型切分模板、prefill/decode worker 资源、batching window、实例队列长度、KV cache 位置、prefix cache 命中情况、session state、链路质量和用户移动趋势。模型切分如果改变 prefill 或 decode 的执行位置，就会改变 batching 的阶段容量和 KV handoff 成本；batching 如果改变请求等待时间和 token 生成节奏，就会影响 request routing 观察到的实例时延和 admission 边界；request/session routing 如果改变流量入口和服务实例选择，又会反过来改变模型实例负载、prefix cache 命中率和后续 context/KV 迁移成本。

因此，三类算法之间不是简单的前后级联关系，而是多时间尺度上的闭环关系。低频部署层决定模型切分、模型副本、SLM/LLM 插件和服务能力布局；中频 serving 层根据 workload 调节 batching、admission、prefill/decode 比例和实例负载均衡；高频 offloading 层根据请求复杂度、用户位置、链路状态和状态位置执行 request routing、workflow placement 和 state migration。协同优化的关键不是把每一层都独立做到局部最优，而是让它们在统一 SLO、资源约束和状态约束下避免互相抵消。

从控制变量看，三类算法的输入和输出也存在明显依赖。模型切分层输出可部署模板、阶段容量、KV handoff cost 和 fallback path；batching 层输出队列等待、active batch occupancy、TTFT/TPOT、KV 水位和阶段 backlog；offloading 层输出请求入口、计算节点、状态处理动作和后续 session 的 KV owner。这些输出又会成为其他层下一轮决策的输入。换言之，协同优化不应只追求一个全局静态最优解，而应维护一个持续更新的系统状态视图。


\phantomsection
\subsection*{7.2 Model / Service Placement -- Request Offloading 协同}
\addcontentsline{toc}{subsection}{7.2 Model / Service Placement -- Request Offloading 协同}

模型/服务部署与请求卸载之间的耦合，是 CEC-LLM 推理系统中最直接的协同优化关系。由于边缘节点无法缓存所有 LLM、SLM、expert、adapter 或工具服务，模型能力是否提前部署会直接决定请求能否在近端被处理，以及是否需要跨边缘或上云。反过来，请求流量、用户会话分布和模型调用频率又会影响下一周期的模型缓存和服务部署收益。

已有工作已经开始显式处理部署规划与请求卸载之间的时间尺度差异。SLM-LLM-Collab \refcite{60} 将 SLM caching 与 local/D2D/edge inference offloading 进行双时间尺度建模，说明模型缓存不是独立部署问题，而是为后续请求卸载提供可行路径。T2DRL \refcite{68} 进一步将 long-context LLM serving 中的 model caching、request management 和资源匹配联系起来，强调上下文长度和移动边缘资源会共同影响模型缓存与请求处理策略。H-MEC-Placement \refcite{71} 虽然不是 LLM-specific 工作，但从 general MEC 角度说明 service placement、task scheduling 和 resource allocation 在规划周期内存在强耦合。LVM-Offload \refcite{72} 则在 foundation model 服务中联合考虑模型选择、推理步数、任务卸载和资源分配。

这一类协同的核心变量包括 model/service placement、SLM/LLM cache、adapter/expert residency、service capacity、request arrival pattern、routing action、cloud fallback 和 cache update period。低频部署层如果只根据历史平均流量放置模型，可能无法应对移动热点和突发请求；高频请求卸载如果只根据当前节点负载选择路径，又可能频繁调用未缓存模型，造成冷启动、跨域传输和质量回退。因此，部署层和卸载层需要形成“长期能力布局 + 短期路径选择 + 在线反馈更新”的闭环。

在 CEC 中，模型部署层应当被视为协同优化框架中的低频规划层，而不是一次性静态配置。其输出不应只是某个模型放在哪个节点，而应包括模型能力布局、服务容量边界、可服务请求类型、切换条件和回退路径。请求卸载层则在给定部署状态下进行高频决策，并把流量、命中率、队列和质量反馈回部署层，用于下一周期更新模型缓存和服务布局。


\phantomsection
\subsection*{7.3 Model Partitioning -- Batching / Scheduling 协同}
\addcontentsline{toc}{subsection}{7.3 Model Partitioning -- Batching / Scheduling 协同}

模型切分与 batching/scheduling 的耦合主要体现在推理阶段的容量匹配和时延控制上。LLM inference 中 prefill、decode、verification 等阶段具有不同资源画像：prefill 更偏计算密集，decode 更偏 memory bandwidth 和低时延 token streaming。模型切分或阶段解耦改变这些阶段所在的资源池，而 batching 和 scheduling 决定请求如何进入这些资源池以及 token 如何被迭代生成。

已有 LLM serving 工作从不同角度揭示了这种耦合。Splitwise 和 DistServe 通过 prefill/decode disaggregation 说明，将两类阶段部署到不同资源池可以提高 goodput，但也会引入阶段间协调和 KV handoff 成本 \refcite{18}, \refcite{7}。Orca \refcite{4} 和 vLLM \refcite{5} 等系统说明，iteration-level scheduling 和 continuous batching 可以显著提高 GPU 利用率，但其效果受到模型部署方式、decode worker 容量和 KV cache 管理策略影响。Llumnix 进一步展示了动态请求迁移、实例调度和 serving 稳定性之间的关系 \refcite{19}。Mooncake 从 KVCache-centric serving 出发，将 KV cache 管理、prefill/decode 分离和请求调度结合起来，说明状态位置会影响 batching 和阶段执行路径 \refcite{21}。Jupiter 将 generative LLM 在 edge devices 上的 collaborative inference 与 pipeline parallelism、speculative decoding 结合起来，进一步说明在 CEC 场景中，阶段拆分和 serving scheduling 必须一起考虑 \refcite{27}。

这一类协同的关键变量包括 partition template、stage boundary、prefill/decode pool capacity、batching window、active sequence set、chunk size、KV transfer time、TTFT/TPOT target 和 worker scaling。模型切分模板不能只描述模型层或阶段放在哪里，还需要暴露阶段容量、KV handoff 成本、batching 可容纳规模和 TTFT/TBT 估计。Batching 层也不能只根据当前队列最大化吞吐，而需要感知切分拓扑、状态位置和用户 SLA。

在 CEC 系统中，模型切分与 batching 的协同尤其需要避免局部优化互相抵消。例如，阶段切分把 prefill 放到区域边缘可以提升 prefill 吞吐，但如果 KV 传回接入边缘的链路不稳定，decode TTFT 和后续 token latency 可能恶化；batching 策略为了追求高吞吐扩大等待窗口，可能让移动用户在 handover 前无法及时完成关键阶段；continuous batching 扩大 active set 能提升利用率，但也会提高 KV 显存压力，使请求迁移和状态复制更困难。因此，协同控制器需要同时观察阶段负载、batch 状态、KV 水位和链路状态。


\phantomsection
\subsection*{7.4 Request / Workflow Routing -- Runtime State Management 协同}
\addcontentsline{toc}{subsection}{7.4 Request / Workflow Routing -- Runtime State Management 协同}

请求路由和 runtime state 管理之间的耦合，是 LLM inference 相比传统 task offloading 最突出的新问题之一。传统 offloading 中，任务完成后状态通常不再影响后续请求；而在 LLM 多轮对话、长上下文和 agent workflow 中，context、KV cache、prefix cache、agent memory 和 tool outputs 会持续影响后续推理。此时，请求应该路由到哪里，不仅取决于哪个节点算力空闲，还取决于相关状态是否已经在该节点或邻近节点存在。

Mooncake、Preble 和 CachedAttention 等 cache-centric serving 工作说明，prefix/KV cache 已经从单纯的内存优化对象变成影响请求调度和服务路径选择的核心资源 \refcite{21}, \refcite{23}, \refcite{24}。FuseSpill \refcite{66} 从显存受限场景出发说明，KV cache spillover 可能导致 decode 请求迁移到辅助 GPU 或其他资源位置，从而改变原有的执行路径。CA-AIGC-Migration \refcite{69} 则直接面向移动边缘场景，将用户 long context 视为可迁移状态，而不是迁移完整模型，从而在回复质量和迁移成本之间做权衡。QoS-LLM-Router \refcite{55} 和 CES-CB \refcite{56} 虽然主要关注 request-level routing，但它们也说明了 routing action 会改变队列状态、batch 组成和后续请求的等待时延。

这一类协同的关键变量包括 session ID、prefix hash、KV block location、state size、state freshness、remote access latency、migration cost、recompute cost、queue state 和 request quality target。对于连续 session，路由动作和状态动作应被联合枚举：在已有 KV 的节点本地执行、将 KV 迁移到新节点执行、远程访问旧节点状态，或在新节点重新 prefill / recompute。不同动作的当前成本和长期收益不同，不能只用单次请求的最小时延判断。

对于 multi-agent workflow，这种耦合更加复杂。MasRouter、Conductor 和 AFLOW 等工作主要从算法层研究如何选择 agent 数量、角色、模型和工作流结构，但在 CEC 系统中，每个 agent 背后的模型、工具、数据库和 memory 可能位于不同物理节点。因此，workflow orchestration 必须进一步转化为 workflow offloading：不仅决定 planner、solver、critic、verifier 等角色如何协作，还要决定这些组件在哪里执行、哪些 context 可以共享、哪些中间状态需要迁移，以及何时调用云端强模型进行 verification。


\phantomsection
\subsection*{7.5 Mobility-aware Hierarchical Control}
\addcontentsline{toc}{subsection}{7.5 Mobility-aware Hierarchical Control}

多用户移动性会同时影响模型部署、batching、请求路由和状态管理，因此它是三类算法协同优化中最重要的外部扰动之一。用户移动会改变区域请求热点，使原有 model placement 失效；会改变接入链路和候选 edge，使 request routing 的当前最优路径失效；会在请求或 workflow 执行过程中触发 handover，使中间状态和后续阶段的执行路径需要重新选择；也会使 long context、KV cache 和 agent memory 留在旧 edge，从而产生状态迁移、远程访问或重算之间的权衡。

已有移动性相关工作提供了局部研究基础。MGCO \refcite{61} 说明，generative computation offloading 不能只根据当前状态短视决策，而需要利用长程移动信息进行规划。DT-MADTO \refcite{65} 通过 digital twin 建模 DAG task、用户移动和中间数据流，说明移动性会改变 workflow/subtask offloading 的传输代价和拥塞状态。CA-AIGC-Migration \refcite{69} 进一步将移动性与 long context migration 结合，说明状态迁移可以成为替代完整服务迁移的轻量机制。T2DRL \refcite{68} 则提示，在 mobile edge 场景中，long-context serving 的模型缓存、请求管理和资源匹配需要与移动场景共同考虑。

这一类协同的关键变量包括 user trajectory、handover probability、future ingress node、service hotspot、state owner、migration bandwidth、session continuation probability 和 workload forecast。低频层根据用户移动和服务热点调整模型/服务布局；中频层根据负载和状态命中率调整 batching 与实例分流；高频层根据 handover、链路抖动和 context/KV 位置进行 request/workflow/state routing。三个时间尺度不需要每次都同时重优化，但必须共享预测和反馈。

现有研究还没有形成完整的 CEC-LLM 跨层闭环。很多工作只优化其中一层：有的关注 request offloading，有的关注 DAG subtask，有的关注 context migration，有的关注 model caching。未来更需要一个 mobility-aware hierarchical controller，使模型能力、请求路径、workflow 结构和状态位置能够随用户移动和 workload 变化自适应调整。


\phantomsection
\subsection*{7.6 算法方案对比与面向华为验证环境的建议}
\addcontentsline{toc}{subsection}{7.6 算法方案对比与面向华为验证环境的建议}

结合第四、五、六章的项目约束，第七章不宜把协同优化写成一个不可落地的全局大优化问题。更合理的定位是：在已有模型切分规划器、业务感知 batching 控制器和长期成本感知请求卸载策略之上，实现一个\textbf{分层状态感知协同控制器}。它不替代底层 runtime 的 continuous batching、PD 分离、KV cache 管理或路由执行逻辑，而是统一读取部署模板、batch 状态、节点负载、链路状态和 KV/session 元数据，在不同时间尺度上修正各层参数。

该控制器的第一部分是统一状态视图。系统需要收集模型/服务部署状态、可用切分模板、prefill/decode worker 容量、batching 参数、队列长度、TTFT/TPOT/E2E latency、KV cache 位置、prefix 命中率、节点显存、链路 RTT/带宽、请求优先级和用户入口变化。只有把这些信息放入同一个状态表，才能判断某个请求到底应该近端执行、跨节点迁移状态、进入更大 batch，还是回退到云端/区域边缘。

第二部分是分层决策周期。低频周期执行模型/服务 placement 与切分模板选择，重点看历史请求分布、未来热点、节点容量和模型能力覆盖；中频周期调节 batching window、batch token budget、prefill/decode ratio、admission threshold 和实例分流；高频周期执行 request routing、workflow placement、KV migrate/recompute/local 动作和异常回退。不同周期之间通过显式接口传递容量、成本和状态，而不是各自维护不一致的局部视图。

第三部分是冲突检测与受控重配置。协同控制器需要识别三类典型冲突：模型切分带来的 KV handoff 成本超过 batching 收益；batching 为吞吐扩大等待窗口导致高优先级请求或移动用户违反 SLO；request routing 因已有 KV 黏附远端节点，长期造成跨节点通信。对于这些冲突，控制器应通过候选模板切换、batch 参数收缩、状态迁移/重算、路由惩罚项和冷却时间进行修正，避免频繁重配置造成调度震荡。

验证设计上，建议第七章对应的实验分成四组。第一组是独立算法组合基线：分别启用第四章切分部署、第五章动态 batching、第六章请求卸载，但不共享状态，只用各自局部观测做决策。第二组是共享状态但无长期预测：三类算法读取统一节点、队列和 KV 元数据，但只优化当前窗口。第三组是分层协同控制：低频更新部署/切分模板，中频调节 batching 参数，高频执行状态感知请求卸载，并把移动入口、KV owner 和未来 session 成本纳入决策。第四组是扰动与消融：改变用户移动比例、链路带宽、KV cache 规模、长 prompt 比例、请求优先级分布和节点显存水位，分别关闭状态共享、移动预测、KV 迁移、batch 参数联动和切分模板联动，观察 P99 TTFT、平均 E2E latency、SLA violation、跨节点通信量、KV 命中率和满足 SLA 的并发数。

因此，第七章面向验证环境最适合落地的是\textbf{统一元数据驱动的分层协同优化框架}。它的算法主线不是一次性求解所有 placement、batching 和 offloading 变量，而是建立“低频部署布局 -- 中频 serving 参数 -- 高频请求/状态路由”的闭环。这样既能和第四、五、六章的单项算法保持边界清晰，也能把协同收益落实到可测指标上：在满足 P99 TTFT SLA 的前提下提升并发能力，降低平均 E2E 时延，减少跨节点状态通信，并提高移动 session 的服务连续性 \refcite{76}。


\phantomsection
\subsection*{7.7 小结}
\addcontentsline{toc}{subsection}{7.7 小结}

总体来看，三类算法的协同优化不是简单地把模型切分、batching 和 task offloading 合并成一个大优化问题，而是要在实际工作揭示的局部耦合关系基础上，建立分层、多时间尺度、状态感知的控制框架。本章将相关耦合关系组织为四类：Model / Service Placement -- Request Offloading 协同、Model Partitioning -- Batching / Scheduling 协同、Request / Workflow Routing -- Runtime State Management 协同，以及 Mobility-aware Hierarchical Control。

已有工作已经分别展示了 model caching 与 inference offloading、prefill/decode disaggregation 与 scheduling、KV cache 与 request routing、context migration 与 mobility-aware serving 之间的耦合，但这些机制仍然缺少面向 CEC-LLM inference 的统一组织。CEC 场景的核心挑战在于，模型能力、batch 状态、请求路径和运行时状态会随着用户移动、负载变化和链路波动持续变化；局部最优策略如果缺少共享状态和长期反馈，容易造成阶段失衡、尾延迟放大、KV 状态黏附或频繁迁移。

面向工程落地，协同优化应采用“统一状态视图 + 分层决策周期 + 受控重配置”的方式推进。低频层负责模型/服务部署和切分模板选择，中频层负责 batching 参数、admission 和阶段队列调节，高频层负责 request/workflow routing 与 context/KV/runtime state 处理。三层通过 TTFT/TPOT、queue length、KV 命中、链路状态、handover 和显存水位等指标持续反馈，从而使模型切分、batching 和 task offloading 不再只是并列算法，而是服务于同一 CEC-LLM 推理目标的协同控制闭环。

